% Appendix D, construction details (restructure plan, part 3). Sources: experiments/protocol.md (§4, §5, §15.5,
% §15.6, §16.5); src/tracecontract/contract.py (requirements_from_proposals, validate_contract, GATE_BLOCKING_WARNINGS),
% src/tracecontract/audit.py, src/tracecontract/verifier.py (the lint); tasks/CS-0*/verifiers/*/spec.json;
% docs/decision_log.md (the D-ids in the text). Moved here: data.tex's controlled sessions and initial state, method.tex's
% contracts, conditions, audit, freezing and label-support verifiers, results_pilot.tex's contract recovery, the
% appendix's verifier and task tables, and the acceptance gate. Every number is a macro (\fv, \plv, \dv).
\section{Construction details}
\label{app:construction}

\subsection{Controlled sessions}
\label{sec:controlled}
Each scenario was hashed and logged locally before its session ran (protocol CS-v1; not a formal preregistration). It
fixes an initial request $U_1$, a correction $U_2$, a neutral reply for clarification questions, and the intended final
requirements with the turn that introduced each (Table~\ref{tab:scenarios}); CS-04's unchanged default is an inferred
requirement without a quote (D41). The scenario writer is the study's agent, the same model as the session agent, and no
person wrote or reviewed a scenario; for \texttt{humanize} the reference operationalisation was written after the agent's
first output had been seen. The session agent, a subagent of an AI coding agent, works in a fresh git copy of a pinned
commit of a permissively licensed Python library, and runs repository code only through wrappers that execute it in an
offline Docker sandbox. Its workspace is captured after $U_1$ (the pre-correction state) and at the end (the reference
solution). The correction arrives through a subagent messaging channel that adds a short harness frame (D10).
Transcripts are audited for access outside the workspace and for reads of sealed material, then sanitised: harness
internals are dropped and e-mail addresses redacted. Each session gives a task family (Table~\ref{tab:tasks}).

\subsection{Contracts}
\label{sec:contracts}
A contract is a set of requirements, each with a statement, a kind (change, invariant, constraint, non-goal), a status
(\emph{explicit}, \emph{inferred}, \emph{superseded} or \emph{unresolved}), the session event that introduced it, and
evidence spans that must match the session text exactly. The same proposer prompt (\texttt{contract:v1}) recovers the
correction-aware contract from the whole session and the initial-only contract from the first user request alone, so
the contracts differ only in what reaches the proposer ($U_1$ alone or the whole session), not in the proposer's
capability; every generator also gets the reference diff (Section~\ref{sec:design-conditions}). As the code implements it (\texttt{src/tracecontract/contract.py}):
\begin{itemize}
\item quotes are aligned to any session event; an explicit requirement that quotes a non-user event is kept, with a
  \texttt{non\_user\_evidence} warning that blocks the acceptance gate;
\item quotes that do not align are dropped with a warning; an explicit requirement left with no aligned quote is
  downgraded to inferred if it carries a rationale, and rejected otherwise;
\item supersession links must point forward in time, active requirements may not conflict, an initial-only contract may
  use no later event, and superseded and unresolved requirements are never scored;
\item inferred requirements, like every requirement of a reference contract, need human review before a task is
  accepted.
\end{itemize}
For a natural task the proposer reads a rendered pull-request thread: the issue alone, or the issue with the review
thread (protocol~\S15.6); a proposal introduced at an event that is not a user turn (the issue, the review comment or a
reply) is rejected.

\subsection{Verifier conditions}
\label{sec:conditions}
Table~\ref{tab:conditions} lists the conditions. A generator writes black-box tests from a contract (the prompt asks for
\dv{design.tests.min} to \dv{design.tests.max} test functions, and the audit's reviser may keep
\dv{design.revise.tests.min} to \dv{design.revise.tests.max}; in the controlled tasks C1, C3 and B2 have
\fv{checks.plain.min} to \fv{checks.plain.max} scored checks and the audited C2 and C4 \fv{checks.audited.min} to
\fv{checks.audited.max}, Table~\ref{tab:verifiers}). Every generator gets
the base repository and the reference diff and runs code only against the base. Protocol~\S4 asks for the same generator
model for every condition, which no record can confirm for the controlled arm (Table~\ref{tab:roles}), and the same
prompt template, apart from the deviations of Appendix~\ref{app:deviations}. The natural arm uses these templates with
marked deviations, checked word by word against their sources and hashed before any natural agent started (D64).

\begin{table}[t]\centering\scriptsize
\caption{The verifier conditions (protocol~\S4, \S15.6). Every generator gets the base repository and the reference diff,
and is told not to over-fit to the reference's implementation choices; plain generators may not execute their tests on
any candidate. The audit grades development candidates only (Section~\ref{sec:design-conditions}).}\label{tab:conditions}
\begin{tabular}{@{}>{\raggedright\arraybackslash}p{0.2\linewidth}>{\raggedright\arraybackslash}p{0.29\linewidth}>{\raggedright\arraybackslash}p{0.16\linewidth}>{\raggedright\arraybackslash}p{0.27\linewidth}@{}}\toprule
Condition & Generator's contract or input & Audit & Notes \\\midrule
C1 initial-only, plain & initial-only contract, recovered from $U_1$ alone & none & single-pass generation \\
C2 initial-only, audited & the same contract & \texttt{audit:v1}, at most \dv{design.audit.revisions.w} revisions &
starts from C1's tests \\
C3 correction-aware, plain & correction-aware contract, recovered from the whole session & none & single-pass
generation \\
C4 correction-aware, audited & the same contract & \texttt{audit:v1}, at most \dv{design.audit.revisions.w} revisions &
starts from C3's tests \\
B1 repository tests & none & none & the repository's tests at the reference state that pass on the reference, the session
agent's added tests among them; candidate code in the reporting process \\
B2 session-LLM baseline & the whole session transcript, without a contract & none & single-pass generation \\
B3 diff-derived (natural arm, secondary) & the initial request and the reference diff, without a contract & none & a test
is kept only if it fails on the base and passes on the reference \\
\bottomrule\end{tabular}\end{table}

C1 and C3 are single-pass; C2 and C4 start from their tests and apply the audit. B1 imports candidate code directly; in
the controlled tasks the reference is the session agent's final state, so B1 includes
\plv{b1.checks.added.min} to \plv{b1.checks.added.max} tests per task that the agent itself added. B2 is a generator
given the whole session transcript and no contract machinery, a competent ``just ask an LLM'' baseline. B3, natural arm
only, gets the initial request and the reference diff; its tests must fail on the base and pass on the reference
(protocol~\S15.6).

\subsection{The audit and freezing}
\label{sec:audit}\label{sec:freezing}
The audit uses development candidates only: the reference, graded \dv{design.audit.reference.repeats.w} times, the
no-op, \dv{design.audit.attacks.w} attacks written from the same contract, each to violate a requirement, and
\dv{design.audit.alternatives.w} valid alternative. It passes when the reference passes every repeat, the no-op and every
attack fail, the alternative passes, every scored requirement maps to a check and a static lint passes; an attack that
fails only checks mapped to other requirements is reported as a finding that does not fail the audit
(\texttt{src/tracecontract/audit.py}). The lint admits imports from the standard library, \texttt{pytest},
\texttt{hypothesis} and the sandbox API only, and rejects parametrised tests, tests whose only calls are to
\texttt{print}, \texttt{len} or \texttt{isinstance}, tests whose assertions are missing or check only a result's
\texttt{ok} flag, and writes to literal paths in the candidate workspace or results directory. In every frozen generated
verifier (C1 to C4, B2) each test calls \texttt{run\_python} directly, and none writes files from the grader process. At
most \dv{design.audit.revisions.w} revisions are allowed, and a verifier whose audit still fails is frozen as it stands,
as CS-01's C2 was (D27). In every pilot audit record, the reference's repeated gradings agreed.

The audit records of CS-01's C2 and CS-03's C4 hide part of what happened (D62). CS-01's C2 record names
\dv{audit.message.named} of the \fv{ref.c2.cs01.failed} checks that fail on the reference, its message cut at
\dv{audit.message.named}. CS-03's C4 record shows no finding after a revision that added \fv{audit.c4.cs03.added} checks,
without showing them. The files are kept as recorded; new audit records keep every failing check and every check a
revision adds or removes.

The pilot's evaluated verifiers (C1 to C4, B1, B2) were frozen by content hash before any generated held-out candidate
existed, except B1 for \texttt{more-itertools}, whose check selection was re-run mechanically after a harness defect had
truncated it, when superseded results using the truncated selection had been seen (D31, D32). Specs first hashed only
their test files; a hash of the scored check list was added to all \dv{pilot.specs.hashed} after the run
\texttt{pilot-v0.3} (D37), and no pilot spec records its freeze time, so the freeze order rests on git history. Every
verifier that grades the fresh attempts, the independent labeler included, was frozen before the first attempt. The
natural verifiers were frozen phase by phase, each before its family's held-out pool (D64, D65).

\subsection{The verifiers behind the labels}
\label{sec:labelers}
The \emph{reference checks} and the \emph{independent labeler} support the labels (Section~\ref{sec:design-labels}) and
are not conditions.
\begin{itemize}
\item The reference checks were written by configuration-A agents from the reference contract, with no recorded read of
  the held-out pool, and validated only on the reference and the no-op.
\item The independent labeler (protocol~\S15.5; D55, D58) is written per family by a configuration-B agent that may read
  only the reference contract, the base snapshot, the harness API, the lint rules and a synthetic test example: no
  reference diff, reference check, verifier, candidate or label. Its checks are per requirement (specified properties,
  invariants checked differentially against the base, quoted examples), and it does not run the repository's tests,
  which D19's regression norm uses.
\item A configuration-A coordinator graded each controlled labeler \dv{fresh.labeler.repeats.w} times on the reference
  and on the development no-op, was to pass back only failing check names, and froze it after at most
  \dv{fresh.labeler.revisions.w} revision. The labelers have \fv{checks.cs01.il}, \fv{checks.cs02.il},
  \fv{checks.cs03.il} and \fv{checks.cs04.il} checks in CS-01 to CS-04 (only CS-03's was revised); each passes the
  reference in every repeat, fails some check of every change requirement on the no-op and covers every active
  requirement. They were frozen from \dv{fresh.labelers.first} to \dv{fresh.labelers.last} UTC on \dv{fresh.date},
  before any fresh attempt, and registered in sidecar files after a first registration that changed the frozen task
  manifests was reverted (D55).
\item A strict audit of the \dv{fresh.labeler.writers.w} writer transcripts, confirmed by a second agent, found no read of
  forbidden material (D58). Its verdicts are \emph{independent} for CS-01, CS-02 and CS-04 and \emph{concerns} for CS-03,
  whose revision prompt also gave the failing checks' kinds, against the names-only rule, and whose reviser tried a
  repository-wide search that was denied; that revision also leaves the attribution of its all-caps checks to
  requirements unreliable, though not its candidate-level verdict.
\item The exploratory agreement of the labelers with D19 on the \plv{heldout} pilot candidates (protocol~\S15.5) has not
  been computed.
\end{itemize}
The natural arm's reference checks go through the review loop of D33 for every family and carry a status record, and
its labelers follow D64's rules (Section~\ref{sec:natural-cases}).

\subsection{Initial state, leakage and the reconstruction check}
\label{sec:initial-state}
\begin{itemize}
\item A controlled base snapshot is the exact tree of the pinned commit (\texttt{git archive}, no history), and the fresh
  attempts start from it; CS-01's base adds a generated version file, recorded as an approximation.
\item The leakage check recorded with each snapshot was vacuous as called: it compared the base with itself (D47).
\item A reconstruction check, added after the pilot, tests each base against its own session from recorded material
  only (D52):
  \begin{itemize}
  \item replaying the session's edits must reproduce the captured pre-correction and final states of every file they
    wrote;
  \item at each file's first edit the replaced text must occur in the base and the new text must not;
  \item file contents that the agent saw while unchanged (Read, Grep, plain \texttt{cat}, \texttt{sed}, \texttt{head})
    must equal the base;
  \item the rebuilt base must equal the manifest recorded before the session, without compiled or backup copies of
    changed files.
  \end{itemize}
  Every controlled base passes, also under the check's second version. In \texttt{humanize}, a test file that the agent
  rewrote with a script before its first edit, and the added version file, of which the session saw part (it matches),
  cannot be fully checked. Since protocol v0.8 every new task must pass a reconstruction check before it enters a run.
\item A natural task's record checks that both commits are reachable, that the base and merge rules hold, that the
  snapshot equals the base tree, and that the reference and the pre-review state reproduce their commits' trees; a
  missing or non-permissive LICENSE rejects the family. NP-03 is approximate (a symbolic link and submodules not
  materialised), and NP-04, NP-05, NP-18, NP-19 and NP-20 are exact (D59; Table~\ref{tab:natural-s10}).
\end{itemize}

\subsection{The acceptance gate}
\label{sec:gate}
Behind the contracts stands a human-review acceptance gate that no task has passed. Human review alone cannot open it,
since the documented procedure lacks what it requires (D51): no command sets an \texttt{eligible} provenance record, and
every record is \texttt{pending\_review}; and the gate checks every contract a task lists, while the initial-only and
correction-aware contracts hold unresolved requirements by design and the blinded review packet covers only the
reference contract (Appendix~\ref{app:annotation}). \texttt{tracecontract accept-task -{}-dry-run} lists every reason.

\subsection{Contract recovery and the tables of the controlled tasks}
\label{sec:recovery}
Contracts recovered from the whole session cover every active requirement of the reference contract in every family.
The initial-only contracts cover from \plv{coverage.initial.min} to \plv{coverage.initial.max} of them
(Table~\ref{tab:coverage}). In \texttt{humanize} the initial-only contract maps to a requirement that its other
requirements contradict (reported apart as ``covered but contradicted''). A second, independent model mapper agrees
closely with the first ($\kappa$ from \plv{kappa.min} to \plv{kappa.max}), but both mappers are the same model, and the
second mapper, and the first for \texttt{tinydb}, received an extra instruction (D41). We corrected a contradiction claim
of the first mapper for \texttt{humanize} after a sandbox check showed that it rested on a wrong assumption about the
base behaviour (\texttt{tasks/CS-01/contracts/mapping\_full\_session.json}, field \texttt{corrections}; D42). Without the
correction, that correction-aware contract would show a contradicting requirement and a covered-but-contradicted
requirement; its coverage would be unchanged. Table~\ref{tab:verifiers} lists the scored checks of every verifier, the
audit revisions of the audited ones, and the checks that the session agent added to the repository's tests;
Table~\ref{tab:tasks} lists the tasks and Table~\ref{tab:scenarios} the scenarios' intended requirements. These
tables, like the pilot's tables of Appendix~\ref{app:pilot}, are frozen with the pilot's run
(\path{scripts/verify_frozen_inputs.py}) and keep the names they were frozen with: there ``full'' and
``full-session'' mean correction-aware (C3, C4), ``initial'' means initial-only (C1, C2), and ``full-session LLM'' is the
session-LLM baseline (B2).

\IfFileExists{generated/coverage_table.tex}{\begin{table}[t]\centering\footnotesize
\caption{Recovered contracts against the reference contract (model mappings, proposals). Coverage: share of active scored reference requirements that a recovered requirement maps to without contradicting it.}\label{tab:coverage}
\begin{tabular}{llrrrrr}\toprule Task & Contract & Coverage & Unsupported & Contradicting & Coverage (mapper 2) & $\kappa$ \\\midrule
CS-03 & full-session & 1.00 (4/4) & 1/5 & 0 & -- & -- \\
CS-03 & initial-only & 0.75 (3/4) & 1/5 & 0 & -- & -- \\
CS-04 & full-session & 1.00 (6/6) & 1/7 & 0 & -- & -- \\
CS-04 & initial-only & 0.50 (3/6) & 0/3 & 0 & -- & -- \\
\bottomrule\end{tabular}\end{table}
}{\unfilled{coverage table}}
\IfFileExists{generated/verifiers_table.tex}{\begin{table}[t]\centering\scriptsize
\caption{Verifiers per task (generated from \texttt{tasks/*/verifiers/*/spec.json}): scored checks. Audited verifiers: audit revisions, whether the audit passed, and ``= plain'' when the audited tests are identical to the plain ones. B1: checks added to the repository's tests by the session agent (not in the base commit). Ref.: reference checks used for labels.}\label{tab:verifiers}
\begin{tabular}{lccccccc}\toprule Task & C1 & C2 & C3 & C4 & B2 & B1 & Ref. \\\midrule
CS-01 & 14 & 14 (1 rev., not passed) & 14 & 14 (0 rev.; = plain) & 14 & 737 (11 agent-added) & 17 \\
CS-02 & 14 & 14 (0 rev.; = plain) & 13 & 13 (1 rev.) & 13 & 239 (12 agent-added) & 28 \\
CS-03 & 14 & 15 (1 rev.) & 12 & 14 (1 rev.) & 13 & 466 (11 agent-added) & 15 \\
CS-04 & 12 & 12 (0 rev.; = plain) & 13 & 13 (0 rev.; = plain) & 14 & 777 (6 agent-added) & 19 \\
\bottomrule\end{tabular}\end{table}
}{\unfilled{verifier table}}
\IfFileExists{generated/tasks_table.tex}{\begin{table}[t]\centering\small
\caption{Tasks (generated from task manifests). Held-out candidates exclude the reference solution and the no-op.}\label{tab:tasks}
\begin{tabular}{lllllr}\toprule Task & Population & Repository & Commit & Verifiers & Held-out \\\midrule
CS-01 & controlled & \texttt{python-humanize/humanize} & \texttt{a950045} & 7 & 3 \\
CS-02 & controlled & \texttt{msiemens/tinydb} & \texttt{edff96a} & 6 & 1 \\
CS-03 & controlled & \texttt{jpvanhal/inflection} & \texttt{88eefaa} & 7 & 8 \\
CS-04 & controlled & \texttt{more-itertools/more-itertools} & \texttt{ceafee8} & 7 & 8 \\
\bottomrule\end{tabular}\end{table}
}{\unfilled{tasks table}}
\IfFileExists{generated/scenario_table.tex}{\begin{table}[t]\centering\footnotesize
\caption{Intended requirements of the registered controlled scenarios, by the user turn that introduced them (generated from \texttt{data/controlled/scenarios}).}\label{tab:scenarios}
\begin{tabular}{llp{0.33\linewidth}p{0.33\linewidth}}\toprule Task & Library & Initial request ($U_1$) & Correction ($U_2$) \\\midrule
CS-01 & \texttt{humanize} & R-default-trim (superseded), R-keep-meaningful & R-opt-in, R-default-unchanged, R-custom-format-unchanged \\
CS-02 & \texttt{tinydb} & R-atomic, R-new-content & R-read-after-write, R-bytes-unchanged, R-no-temp-files, R-original-intact \\
CS-03 & \texttt{inflection} & R-acronym & R-all-caps-input, R-humanize-unchanged, R-underscore-unchanged \\
CS-04 & \texttt{more-itertools} & R-step & R-negative-step, R-zero-step, R-positional-ordering, R-lazy, R-default-unchanged \\
\bottomrule\end{tabular}\end{table}
}{\unfilled{scenario table}}
